\chapter{Functions}
\index{functions}
\label{ch:functions}

Functions are the primary organizational unit in Scheme. There are no classes and no methods---programs are built by defining and composing functions, with libraries (Chapter~\ref{ch:libraries}) grouping them into reusable units. This chapter covers defining functions, recursion, closures, tail-call optimization, variadic arguments and \texttt{case-lambda}, and returning multiple values.

\section{Defining Functions}
\index{define}

You have seen the shorthand form for defining functions:

\begin{lstlisting}[style=repl]
> (define (square x) (* x x))
> (square 7)
49
\end{lstlisting}

This is syntactic sugar for binding a name to a \texttt{lambda}\index{lambda} expression:

\begin{lstlisting}[style=repl]
> (define square (lambda (x) (* x x)))
> (square 7)
49
\end{lstlisting}

Both forms are equivalent. The shorthand \texttt{(define (name params) body)} is more common in practice, but understanding the \texttt{lambda} form is essential because anonymous functions appear everywhere in Scheme.

\subsection{Multiple Expressions in the Body}

Unlike Python's \texttt{lambda} (limited to one expression), a Scheme lambda body can contain multiple expressions. The value of the last expression is returned:

\begin{lstlisting}[style=repl]
> (define (greet name)
    (display "Hello, ")
    (display name)
    (newline)
    (string-append "Greeted " name))
> (greet "Alice")
Hello, Alice
"Greeted Alice"
\end{lstlisting}

The first three expressions produce side effects (printing). The last expression's value is the function's return value.

\subsection{Functions with Multiple Parameters}

\begin{lstlisting}[style=repl]
> (define (hypotenuse a b)
    (sqrt (+ (* a a) (* b b))))
> (hypotenuse 3 4)
5
\end{lstlisting}

\section{Anonymous Functions}
\index{lambda}\index{anonymous functions}

You do not need to name every function. \texttt{lambda} creates a function value directly:

\begin{lstlisting}[style=repl]
> ((lambda (x) (* x x)) 5)
25
> (map (lambda (x) (+ x 10)) '(1 2 3 4))
(11 12 13 14)
\end{lstlisting}

Anonymous functions are particularly useful when passing short operations to higher-order functions like \texttt{map} and \texttt{filter}. You will use them constantly.

\begin{note}[Coming from Python]
\texttt{(lambda (x) (* x x))} is Scheme's equivalent of Python's \texttt{lambda x: x * x}. The difference is that Scheme's version can contain multiple expressions and arbitrary complexity---there is no need for a separate \texttt{def} just because the function has more than one line.
\end{note}

\section{Recursion}
\index{recursion}

Scheme has no dedicated \texttt{for} or \texttt{while} statement---aside from the \texttt{do} form from Chapter~\ref{ch:control-flow}, repetition is expressed through recursion, and the language is designed to make recursion efficient and natural.

\subsection{Simple Recursion}

The classic factorial:

\begin{lstlisting}[style=repl]
> (define (factorial n)
    (if (= n 0)
        1
        (* n (factorial (- n 1)))))
> (factorial 5)
120
> (factorial 20)
2432902008176640000
\end{lstlisting}

The structure is always the same: a base case that returns a known value, and a recursive case that reduces the problem.

Here is a function that computes the nth Fibonacci number:

\begin{lstlisting}[style=repl]
> (define (fib n)
    (cond
      ((= n 0) 0)
      ((= n 1) 1)
      (else (+ (fib (- n 1))
               (fib (- n 2))))))
> (fib 10)
55
\end{lstlisting}

\begin{warning}[Exponential recursion]
This naive Fibonacci is O($2^n$)---it recomputes the same subproblems repeatedly. For large \texttt{n}, use the tail-recursive version shown later in this chapter.
\end{warning}

\subsection{Recursion on Lists}

Chapter~\ref{ch:lists} showed that a list is either empty or an element followed by a smaller list. That structure dictates the standard recursive pattern: check for empty (base case), process the first element, recurse on the rest. Counting elements:

\begin{lstlisting}
(define (my-length lst)
  (if (null? lst)
      0
      (+ 1 (my-length (cdr lst)))))
\end{lstlisting}

\begin{lstlisting}[style=repl]
> (my-length '(a b c d))
4
\end{lstlisting}

Summing a list of numbers:

\begin{lstlisting}
(define (sum-list lst)
  (if (null? lst)
      0
      (+ (car lst) (sum-list (cdr lst)))))
\end{lstlisting}

\begin{lstlisting}[style=repl]
> (sum-list '(1 2 3 4 5))
15
\end{lstlisting}

The same shape transforms a list into a new one:

\begin{lstlisting}[style=repl]
> (define (my-map f lst)
    (if (null? lst)
        '()
        (cons (f (car lst))
              (my-map f (cdr lst)))))
> (my-map square '(1 2 3 4 5))  ; square was defined above
(1 4 9 16 25)
\end{lstlisting}

\begin{note}[Coming from Python]
This recursive pattern replaces Python's \texttt{for} loop. Instead of iterating with a counter, you decompose the list into head and tail. This style takes practice, but it eliminates off-by-one errors and makes the flow of data explicit.
\end{note}

\subsection{Recursion on Nested Lists}

When lists nest, the recursion branches: if the first element is itself a pair, recurse into it as well as into the rest. Here is a function that flattens a nested list into a single level:

\begin{lstlisting}
(define (flatten lst)
  (cond
    ((null? lst) '())
    ((pair? (car lst))
     (append (flatten (car lst))
             (flatten (cdr lst))))
    (else
     (cons (car lst) (flatten (cdr lst))))))
\end{lstlisting}

\begin{lstlisting}[style=repl]
> (flatten '(1 (2 3) (4 (5 6)) 7))
(1 2 3 4 5 6 7)
\end{lstlisting}

\section{Tail Calls and Tail Recursion}
\index{tail call optimization}\index{TCO}

When a recursive call is the \emph{last thing} a function does---the final expression with nothing left to do after it returns---it is in \emph{tail position}\index{tail position}. Kaappi (and all conforming Schemes) guarantee that tail calls reuse the current stack frame instead of allocating a new one.

This means tail-recursive functions use constant stack space, just like a loop.

\subsection{Non-Tail vs.\ Tail Recursive}

The factorial above is \emph{not} tail-recursive because after the recursive call, it still needs to multiply:

\begin{lstlisting}
;; NOT tail-recursive: (* n ...) wraps the recursive call
(define (factorial n)
  (if (= n 0) 1
      (* n (factorial (- n 1)))))
\end{lstlisting}

Here is the tail-recursive version using an accumulator:

\begin{lstlisting}[style=repl]
> (define (factorial n)
    (define (loop i acc)
      (if (= i 0)
          acc
          (loop (- i 1) (* acc i))))
    (loop n 1))
> (factorial 20)
2432902008176640000
\end{lstlisting}

The call to \texttt{loop} is the last expression---nothing wraps it---so it runs without growing the stack, even for \texttt{(factorial 1000000)}. (The inner \texttt{define} creates a helper procedure visible only inside \texttt{factorial}; these \emph{internal definitions} are covered in Chapter~\ref{ch:bindings}.)

\subsection{Named \texttt{let} for Loops}
\index{named let}

Named \texttt{let} is the idiomatic way to write a tail-recursive loop. It defines a local function and calls it immediately:

\begin{lstlisting}[style=repl]
> (let loop ((i 1))
    (when (<= i 5)
      (display i)
      (display " ")
      (loop (+ i 1))))
1 2 3 4 5
\end{lstlisting}

\texttt{loop} is a name you choose---it is bound to a function with parameter \texttt{i}. Calling \texttt{(loop (+ i 1))} recurses with the next value. Because the call is in tail position, it runs in constant space. Chapter~\ref{ch:bindings} returns to named \texttt{let} and shows how it relates to the other \texttt{let} forms.

Here is a tail-recursive Fibonacci using named \texttt{let}:

\begin{lstlisting}[style=repl]
> (define (fib n)
    (let loop ((i n) (a 0) (b 1))
      (if (= i 0)
          a
          (loop (- i 1) b (+ a b)))))
> (fib 50)
12586269025
\end{lstlisting}

This version is O(n), computing \texttt{(fib 50)} instantly.

\begin{note}[Coming from Python]
Named \texttt{let} is Scheme's equivalent of a \texttt{while} loop with loop variables. The ``loop variables'' are the parameters, and each ``iteration'' is a tail call with updated values. Once you internalize this pattern, you will find it as natural as \texttt{for i in range(...)}.
\end{note}

\section{Closures}
\index{closures}

A \emph{closure} is a function that remembers the environment where it was created. Variables from the enclosing scope remain accessible even after that scope has exited.

\begin{lstlisting}[style=repl]
> (define (make-adder n)
    (lambda (x) (+ x n)))
> (define add5 (make-adder 5))
> (define add10 (make-adder 10))
> (add5 3)
8
> (add10 3)
13
\end{lstlisting}

Each call to \texttt{make-adder} creates a new closure that captures its own copy of \texttt{n}. The two closures \texttt{add5} and \texttt{add10} are independent.

\subsection{Closures with Mutable State}

Closures can capture mutable variables, creating private state. The \texttt{(let ((n 0)) ...)} form below introduces a local variable \texttt{n} bound to 0---Chapter~\ref{ch:bindings} covers \texttt{let} fully; for now, read it as ``create a variable \texttt{n} visible only inside this body'':

\begin{lstlisting}[style=repl]
> (define (make-counter)
    (let ((n 0))
      (lambda ()
        (set! n (+ n 1))
        n)))
> (define c1 (make-counter))
> (define c2 (make-counter))
> (c1)
1
> (c1)
2
> (c1)
3
> (c2)
1
\end{lstlisting}

Each counter has its own private \texttt{n}. This is how you encapsulate state in Scheme without classes.

\begin{note}[Coming from Python]
This is the same concept as Python closures. The difference is that Scheme uses closures as a primary organizational tool (where Python might reach for a class), and \texttt{set!}\index{set!} makes the captured variable mutable.
\end{note}

\section{Variadic Functions}
\index{variadic functions}\index{rest arguments}

Functions can accept a variable number of arguments using \emph{dot notation} in the parameter list:

\begin{lstlisting}[style=repl]
> (define (sum . numbers)
    (let loop ((lst numbers) (acc 0))
      (if (null? lst)
          acc
          (loop (cdr lst) (+ acc (car lst))))))
> (sum 1 2 3)
6
> (sum 1 2 3 4 5)
15
\end{lstlisting}

The dot before \texttt{numbers} means ``collect all remaining arguments into a list.'' You can also have required parameters before the rest:

\begin{lstlisting}[style=repl]
> (define (format-message prefix . parts)
    (apply string-append prefix ": " parts))
> (format-message "ERROR" "file not found")
"ERROR: file not found"
\end{lstlisting}

\section{\texttt{case-lambda}: Optional Arguments}
\index{case-lambda}\index{optional arguments}

Python expresses optional arguments with default values: \texttt{def greet(name, greeting="Hello")}. Scheme's answer is \texttt{case-lambda}, which builds a procedure from several clauses, one per parameter list. When the procedure is called, the clause matching the number of arguments runs:

\begin{lstlisting}[style=repl]
> (define greet
    (case-lambda
      ((name) (greet name "Hello"))
      ((name greeting)
       (string-append greeting ", " name "!"))))
> (greet "Alice")
"Hello, Alice!"
> (greet "Alice" "Howdy")
"Howdy, Alice!"
\end{lstlisting}

The one-argument clause supplies the default by calling the two-argument clause. A clause's parameter list can also use dot notation---\texttt{((x . rest) ...)}---to catch every remaining arity, just as in a variadic \texttt{define}.

\section{\texttt{apply}}
\index{apply}

\texttt{apply} calls a function with arguments from a list:

\begin{lstlisting}[style=repl]
> (apply + '(1 2 3 4 5))
15
> (apply max '(3 7 2 9 4))
9
> (apply string-append '("hello" " " "world"))
"hello world"
\end{lstlisting}

This is essential when you have a list of arguments and need to pass them to a function that expects individual parameters.

\section{Multiple Return Values}
\index{values}\index{multiple values}

Scheme functions can return multiple values using \texttt{values}:

\begin{lstlisting}[style=repl]
> (define (divide a b)
    (values (quotient a b) (remainder a b)))
\end{lstlisting}

The caller must explicitly receive the values. The fundamental receiver is \texttt{call-with-values}\index{call-with-values}, which hands them to another function as its arguments:

\begin{lstlisting}[style=repl]
> (call-with-values (lambda () (divide 17 5)) list)
(3 2)
\end{lstlisting}

In practice you will usually reach for the more convenient \texttt{let-values}\index{let-values} (covered in detail in Chapter~\ref{ch:bindings}), which binds each value to a name:

\begin{lstlisting}[style=repl]
> (let-values (((q r) (divide 17 5)))
    (display "Quotient: ") (display q) (newline)
    (display "Remainder: ") (display r) (newline))
Quotient: 3
Remainder: 2
\end{lstlisting}

\begin{note}[Coming from Python]
This is similar to Python's tuple unpacking: \texttt{q, r = divmod(17, 5)}. Scheme's version is slightly more verbose but achieves the same result.
\end{note}

\section{Function Composition Patterns}

A few patterns come up frequently when working with functions.

\subsection{Functions that Return Functions}
\index{compose}

\begin{lstlisting}[style=repl]
> (define (compose f g)
    (lambda (x) (f (g x))))
> (define inc-then-square (compose square (lambda (x) (+ x 1))))
> (inc-then-square 4)
25
\end{lstlisting}

\subsection{Partial Application}
\index{partial application}

\begin{lstlisting}[style=repl]
> (define (partial f . fixed-args)
    (lambda rest-args
      (apply f (append fixed-args rest-args))))
> (define double (partial * 2))
> (double 5)
10
> (define add1 (partial + 1))
> (add1 99)
100
\end{lstlisting}

These patterns---composition, partial application, and higher-order functions in general---are the tools Scheme gives you in place of classes and inheritance. They are simpler, more composable, and often more powerful.

\section*{Summary}

\begin{itemize}
    \item Define named functions with \texttt{define} and anonymous ones with \texttt{lambda}; both create the same kind of procedure.
    \item Recursion is Scheme's primary looping mechanism, and tail-recursive calls run in constant space.
    \item Named \texttt{let} expresses loops as local tail recursion.
    \item Closures capture their defining environment, enabling private, mutable state.
    \item Variadic functions use a rest parameter, \texttt{case-lambda} dispatches on argument count (Scheme's optional arguments), and \texttt{apply} spreads a list into individual arguments.
    \item Procedures can return several results with \texttt{values}, received with \texttt{call-with-values} or \texttt{let-values}.
\end{itemize}

\section*{Exercises}

\begin{exercise}[Tail-Recursive Summation]
Write a function \code{sum-to} that computes the sum $1 + 2 + \cdots + n$ using tail recursion with an accumulator. Verify that it handles large values without stack overflow:

\begin{lstlisting}[style=repl]
> (sum-to 10)
55
> (sum-to 1000000)
500000500000
\end{lstlisting}

Compare with the formula $n(n+1)/2$ to check correctness.
\end{exercise}

\begin{exercise}[Accumulator Closure]
Write a function \code{make-accumulator} that returns a closure. Each time the closure is called with a number, it adds that number to a running total and returns the new total:

\begin{lstlisting}[style=repl]
> (define acc (make-accumulator))
> (acc 10)
10
> (acc 25)
35
> (acc 5)
40
\end{lstlisting}

In Python, you would use a class with \texttt{\_\_call\_\_} or a closure over a \texttt{nonlocal} variable. In Scheme, \code{set!} and a \code{let}-bound variable do the job.
\end{exercise}

\begin{exercise}[Left-to-Right Composition]
The chapter defined \code{compose}, which applies the \emph{rightmost} function first: \code{((compose f g) x)} computes \code{(f (g x))}. Many pipelines read more naturally the other way around. Write \code{pipe}, which applies its functions left to right---\code{((pipe f g) x)} computes \code{(g (f x))}:

\begin{lstlisting}[style=repl]
> (define inc-then-double
    (pipe (lambda (x) (+ x 1))
          (lambda (x) (* x 2))))
> (inc-then-double 3)
8
> (define exclaim
    (pipe (lambda (s) (string-append s "!"))
          string-upcase))
> (exclaim "hello")
"HELLO!"
\end{lstlisting}

Can you define \code{pipe} in terms of \code{compose}? Bonus (after reading Chapter~\ref{ch:hof}): extend \code{pipe} to accept any number of functions using rest arguments and \code{fold}.
\end{exercise}

\begin{exercise}[Variadic Maximum]
Write a function \code{my-max} that takes one or more numeric arguments and returns the largest. Use a required first parameter plus rest arguments:

\begin{lstlisting}[style=repl]
> (my-max 3)
3
> (my-max 3 7 2 9 4)
9
> (my-max -5 -1 -10)
-1
\end{lstlisting}

The built-in \code{max} already does this, so name yours \code{my-max}. Use a named \code{let} to walk the rest arguments.
\end{exercise}

\begin{exercise}[Tail-Recursive Fibonacci]
The naive Fibonacci function shown earlier in this chapter is O($2^n$). Rewrite it as a tail-recursive function using two accumulators \code{a} and \code{b}:

\begin{lstlisting}[style=repl]
> (fib 0)
0
> (fib 1)
1
> (fib 10)
55
> (fib 50)
12586269025
\end{lstlisting}

Hint: at each step, the new \code{a} becomes the old \code{b}, and the new \code{b} becomes \code{a + b}. This mirrors the named \code{let} version from the tail-call section---try writing it from scratch without looking back.
\end{exercise}

\begin{exercise}[Reversing a List]
Write a function \code{my-reverse} that reverses a list using recursion. Use an accumulator parameter to make it tail-recursive:

\begin{lstlisting}[style=repl]
> (my-reverse '(1 2 3 4 5))
(5 4 3 2 1)
> (my-reverse '())
()
> (my-reverse '(a))
(a)
\end{lstlisting}

In Python, this would be \texttt{lst[::-1]}. Your Scheme version should be O(n) in both time and space.
\end{exercise}

\begin{exercise}[Appending Two Lists]
Write \code{my-append} that concatenates two lists without using the built-in \code{append}:

\begin{lstlisting}[style=repl]
> (my-append '(1 2 3) '(4 5 6))
(1 2 3 4 5 6)
> (my-append '() '(a b))
(a b)
> (my-append '(x) '())
(x)
\end{lstlisting}

Hint: recurse on the first list. What is the base case? Why does this function copy the first list but share the second?
\end{exercise}

\begin{exercise}[Deep Reverse]
Write a function \code{deep-reverse} that reverses a list at every level of nesting. Unlike your plain \code{my-reverse} from earlier, this function must also reverse any sublists it encounters:

\begin{lstlisting}[style=repl]
> (deep-reverse '(1 2 3))
(3 2 1)
> (deep-reverse '(1 (2 3) (4 (5 6)) 7))
(7 ((6 5) 4) (3 2) 1)
> (deep-reverse '((a b) (c d)))
((d c) (b a))
> (deep-reverse '())
()
\end{lstlisting}

This requires two recursive cases: when the \code{car} is itself a pair (deep-reverse it before collecting), and when it is an atom (keep it as-is). Branch on \code{(pair? (car lst))} the way the \code{flatten} example from this chapter does, but collect results into an accumulator so each level comes out reversed.
\end{exercise}

The next chapter covers local bindings in detail---including a closer look at named \texttt{let}, which you met here as a looping construct.

\chapterend
